\hspace{3mm}

\centerline{\bf \Huge Домашнее задание. Турниры}

\hspace{3mm}

\problem{1} Петя, Саша и Миша играют в теннис на вылет. Игра на вылет означает, что в каждой партии играют двое, а третий ждёт. Проигравший партию уступает место третьему и в следующей партии сам становится ждущим. Петя сыграл всего 12 партий, Саша — 7 партий, Миша — 11 партий. Сколько раз Петя выиграл у Саши?

\hspace{3mm}

\problem{2} Команды провели турнир по футболу в один круг (каждая с каждой сыграла один раз, победа - 3 очка, ничья - 1, поражение - 0). Оказалось, что единоличный победитель набрал менее 50\% от количества очков, возможного для одного участника. Какое наименьшее количество команд могло участвовать в турнире?

\hspace{3mm}

\problem{3} В гандбольном турнире в один круг (победа - 2 очка, ничья - 1 очко, поражение - 0) приняло участие 16 команд. Все команды набрали разное количество очков, причем команда, занявшая 7 место, набрала 21 очко. Докажите, что победившая команда хотя бы один раз сыграла вничью.